\documentclass[11pt, a4paper]{article}

\usepackage[T1]{fontenc}
\usepackage[utf8]{inputenc}
\usepackage{tgpagella}
\usepackage[spanish, es-noshorthands]{babel}
\usepackage{booktabs}
\usepackage[most]{tcolorbox}
\usepackage[margin=2.5cm]{geometry}
\usepackage{fancyhdr}

\pagestyle{fancy}
\fancyhf{}
\setlength{\headheight}{16pt}
\lhead{\textbf{Ing. Mecatrónica}}
\chead{Notas de Pre-Lab / Lab 3}
\rhead{Circuitos Electrónicos II}
\renewcommand{\headrulewidth}{0.4pt}
\definecolor{ucbblue}{RGB}{0, 51, 102}

\begin{document}

\begin{center}
    \Large\textbf{\textcolor{ucbblue}{Notas Docentes: Pre-Lab y Ejecución Lab 3}}[cite: 10]
\end{center}

\section*{Evidencia Ideal Esperada del Estudiante[cite: 10]}
\begin{itemize}
    \item Derivación simbólica del amplificador diferencial empleando KCL y divisor de tensión[cite: 10].
    \item Verificación matemática explícita de los ratios resistivos ($\frac{R_2}{R_1} = \frac{R_4}{R_3}$)[cite: 10].
    \item Cálculo numérico del resultado ideal para cada condición elegida[cite: 10].
    \item Comparación $V_{o,\mathrm{ideal}}$ vs LTSpice $V_{o,\mathrm{sim}}$[cite: 10].
\end{itemize}

\section*{Errores Conceptuales Comunes a Vigilar[cite: 10]}
\begin{itemize}
    \item Confundir $V_{OS}$ interno con una diferencia de tensión aplicada externamente[cite: 10].
    \item Tratar la aproximación $V_+ \approx V_-$ como una verdad física exacta, ignorando la ganancia finita y el offset[cite: 10].
    \item Asumir que $I_B = 0$ físicamente, ignorando las caídas de tensión en las resistencias de fuente[cite: 10].
    \item Atribuir automáticamente todo error al Op-Amp antes de auditar el cableado[cite: 10].
    \item Confundir mismatch resistivo externo con offset interno del circuito integrado[cite: 10].
    \item Creer que un error nulo en LTSpice constituye validación del circuito físico[cite: 10].
\end{itemize}

\section*{Errores de Simulación Frecuentes en el Pre-Lab[cite: 10]}
\begin{itemize}
    \item Utilizar un modelo distinto del dispositivo planificado sin registrar la divergencia[cite: 10].
    \item Omitir las fuentes de alimentación DC requeridas por el macromodelo[cite: 10].
    \item Invertir las entradas (conectar la señal a $V_+$ en lugar de $V_-$ o viceversa)[cite: 10].
    \item Usar una referencia de tierra en SPICE que difiera topológicamente de la experimental[cite: 10].
    \item Asumir ciegamente que el macromodelo implementa absolutamente todas las no idealidades del datasheet[cite: 10].
\end{itemize}

\section*{Framework de Troubleshooting (Semana 8)[cite: 10]}
Obligar a los estudiantes a descartar causas en este orden de menor a mayor abstracción:
\begin{enumerate}
    \item \textbf{Supply rails:} ¿Están presentes? ¿Polaridad correcta? ¿Nivel compatible?[cite: 10]
    \item \textbf{Ground / reference:} ¿Se mide con la referencia prevista? ¿Hay tierras flotantes?[cite: 10]
    \item \textbf{Pinout:} ¿Entradas, salidas y alimentación coinciden con el datasheet?[cite: 10]
    \item \textbf{Feedback:} ¿Existe realimentación negativa y al pin correcto?[cite: 10]
    \item \textbf{Resistor values:} Medir los valores reales y comprobar las razones (mismatch)[cite: 10].
    \item \textbf{Input signals:} Medir $V_1$ y $V_2$ en los pines, no asumir que la fuente entrega el valor exacto del display[cite: 10].
    \item \textbf{Measurement setup:} Verificar rango, canal y conexión del instrumento[cite: 10].
    \item \textbf{Real-device nonideality:} Solo después de descartar lo anterior, invocar $V_{OS}, I_B, I_{OS}$[cite: 10].
\end{enumerate}

\end{document}
